\documentclass[10pt,a4paper]{amsart}
\usepackage{amsmath,amssymb,amsthm,graphicx}
\usepackage{verbatim}
\newenvironment{reference}{\comment}{\endcomment}
%\newenvironment{reference}{}{}
\newenvironment{slogan}{\comment}{\endcomment}
\newenvironment{history}{\comment}{\endcomment}
% For commutative diagrams we use Xy-pic
\usepackage[all]{xy}

% We use 2cell for 2-commutative diagrams.
\xyoption{2cell}
% Theorem environments.
%
\theoremstyle{plain}
\newtheorem{theorem}[subsection]{Theorem}
\newtheorem{proposition}[subsection]{Proposition}
\newtheorem{lemma}[subsection]{Lemma}

\theoremstyle{definition}
\newtheorem{definition}[subsection]{Definition}
\newtheorem{example}[subsection]{Example}
\newtheorem{exercise}[subsection]{Exercise}
\newtheorem{situation}[subsection]{Situation}

\theoremstyle{remark}
\newtheorem{remark}[subsection]{Remark}
\newtheorem{remarks}[subsection]{Remarks}

\numberwithin{equation}{subsection}

% Macros
%
\def\lim{\mathop{\mathrm{lim}}\nolimits}
\def\colim{\mathop{\mathrm{colim}}\nolimits}
\def\Spec{\mathop{\mathrm{Spec}}}
\def\Hom{\mathop{\mathrm{Hom}}\nolimits}
\def\Ext{\mathop{\mathrm{Ext}}\nolimits}
\def\SheafHom{\mathop{\mathcal{H}\!\mathit{om}}\nolimits}
\def\SheafExt{\mathop{\mathcal{E}\!\mathit{xt}}\nolimits}
\def\Sch{\mathit{Sch}}
\def\Mor{\mathop{\mathrm{Mor}}\nolimits}
\def\Ob{\mathop{\mathrm{Ob}}\nolimits}
\def\Sh{\mathop{\mathit{Sh}}\nolimits}
\def\NL{\mathop{N\!L}\nolimits}
\def\CH{\mathop{\mathrm{CH}}\nolimits}
\def\proetale{{pro\text{-}\acute{e}tale}}
\def\etale{{\acute{e}tale}}
\def\QCoh{\mathit{QCoh}}
\def\Ker{\mathop{\mathrm{Ker}}}
\def\Im{\mathop{\mathrm{Im}}}
\def\Coker{\mathop{\mathrm{Coker}}}
\def\Coim{\mathop{\mathrm{Coim}}}

% Boxtimes
%
\DeclareMathSymbol{\boxtimes}{\mathbin}{AMSa}{"02}

%
% Macros for moduli stacks/spaces
%
\def\QCohstack{\mathcal{QC}\!\mathit{oh}}
\def\Cohstack{\mathcal{C}\!\mathit{oh}}
\def\Spacesstack{\mathcal{S}\!\mathit{paces}}
\def\Quotfunctor{\mathrm{Quot}}
\def\Hilbfunctor{\mathrm{Hilb}}
\def\Curvesstack{\mathcal{C}\!\mathit{urves}}
\def\Polarizedstack{\mathcal{P}\!\mathit{olarized}}
\def\Complexesstack{\mathcal{C}\!\mathit{omplexes}}
% \Pic is the operator that assigns to X its picard group, usage \Pic(X)
% \Picardstack_{X/B} denotes the Picard stack of X over B
% \Picardfunctor_{X/B} denotes the Picard functor of X over B
\def\Pic{\mathop{\mathrm{Pic}}\nolimits}
\def\Picardstack{\mathcal{P}\!\mathit{ic}}
\def\Picardfunctor{\mathrm{Pic}}
\def\Deformationcategory{\mathcal{D}\!\mathit{ef}}
\usepackage{cleveref}
\crefname{originalreference}{original reference}{original references}
\Crefname{originalreference}{Original reference}{Original references}
\makeatletter
\newlabel{decent-spaces-lemma-locally-separated-decent}{{decent-spaces-lemma-locally-separated-decent}{source}{}{sourcepointer-1}{}}
\newlabel{decent-spaces-lemma-locally-separated-decent@cref}{{[originalreference][0][]decent-spaces-lemma-locally-separated-decent}{[1][0][]source}{}{}{}}
\makeatother
\crefname{originalreference}{original reference}{original references}
\Crefname{originalreference}{Original reference}{Original references}
\makeatletter
\newlabel{decent-spaces-section-reasonable-decent}{{decent-spaces-section-reasonable-decent}{source}{}{sourcepointer-2}{}}
\newlabel{decent-spaces-section-reasonable-decent@cref}{{[originalreference][0][]decent-spaces-section-reasonable-decent}{[1][0][]source}{}{}{}}
\makeatother
\crefname{originallemma}{original lemma}{original lemmas}
\Crefname{originallemma}{Original lemma}{Original lemmas}
\makeatletter
\newlabel{lemma-diagonal-diagonal}{{lemma-diagonal-diagonal}{source}{}{sourcepointer-3}{}}
\newlabel{lemma-diagonal-diagonal@cref}{{[originallemma][0][]lemma-diagonal-diagonal}{[1][0][]source}{}{}{}}
\makeatother
\crefname{originalreference}{original reference}{original references}
\Crefname{originalreference}{Original reference}{Original references}
\makeatletter
\newlabel{spaces-more-groupoids-lemma-group-scheme-over-field-separated}{{spaces-more-groupoids-lemma-group-scheme-over-field-separated}{source}{}{sourcepointer-4}{}}
\newlabel{spaces-more-groupoids-lemma-group-scheme-over-field-separated@cref}{{[originalreference][0][]spaces-more-groupoids-lemma-group-scheme-over-field-separated}{[1][0][]source}{}{}{}}
\makeatother
\crefname{originalreference}{original reference}{original references}
\Crefname{originalreference}{Original reference}{Original references}
\makeatletter
\newlabel{spaces-more-groupoids-proposition-group-space-scheme-over-field}{{spaces-more-groupoids-proposition-group-space-scheme-over-field}{source}{}{sourcepointer-5}{}}
\newlabel{spaces-more-groupoids-proposition-group-space-scheme-over-field@cref}{{[originalreference][0][]spaces-more-groupoids-proposition-group-space-scheme-over-field}{[1][0][]source}{}{}{}}
\makeatother
\crefname{originalreference}{original reference}{original references}
\Crefname{originalreference}{Original reference}{Original references}
\makeatletter
\newlabel{spaces-over-fields-lemma-locally-quasi-finite-over-field}{{spaces-over-fields-lemma-locally-quasi-finite-over-field}{source}{}{sourcepointer-6}{}}
\newlabel{spaces-over-fields-lemma-locally-quasi-finite-over-field@cref}{{[originalreference][0][]spaces-over-fields-lemma-locally-quasi-finite-over-field}{[1][0][]source}{}{}{}}
\makeatother
\title{Six conditions under which the automorphism group of a point is a scheme}
\author{The Stacks Project Authors}
\date{stacks-project:stacks-morphisms.tex:local-snapshot-20261012; GNU FDL 1.2 or later}
\begin{document}
\maketitle
\noindent\textit{Editorial note.} All delivered mathematical text is the author's own, copied from the pinned Stacks Project chapter source sources/stacks-algebraic-stacks-morphisms/pinned-original-source.tex; no statement or proof step was rewritten, reordered or removed.\par\medskip
\noindent\textit{Editorial note.} Preamble provenance: the Stacks Project keeps one single shared preamble for every chapter; this card takes it directly from the staged sources/stacks-algebraic-stacks-morphisms/preamble.tex as three verbatim spans (lines 8-12, 14-18 and 154-217), so no derived file was written and nothing was added or redefined. The wrapper class and the base packages are supplied by the assembler.\par\medskip
\noindent\textit{Editorial note.} Bibliography provenance: neither the counted statement nor its proof contains a citation (this lane rescans the whole span before writing the card), so the card carries no bibliography and no entry was reconstructed.\par\medskip
\noindent\textit{Editorial note.} Selection provenance: this card was opened under the file-based control contract (CONTROL.md), with the mechanical pre-card gate run on its statement and proof interval before the card was written.\par\medskip
\noindent\textit{Allowed standard prerequisites (editorial).}\begin{itemize}
\item The equivalences between separatedness conditions on the diagonal and on the inertia. Invoked at the start of the counted proof (source lines 3957-3960). Reduces three of the six conditions to the remaining three.
\item The statement that a locally separated algebraic space over a field is decent, and that a decent group space over a field is separated. Invoked inside the counted proof (source lines 3961-3970) as cross-chapter pointers. Upgrades quasi-separatedness or local separatedness to separatedness.
\item The statement that a group space over a field which is separated is a scheme. Invoked at the end of the counted proof (source lines 3968-3971) as a cross-chapter pointer. Gives the conclusion in the first two cases.
\item The statement that a group space locally quasi-finite over a field is a scheme. Invoked at the end of the counted proof (source lines 3972-3974) as a cross-chapter pointer. Gives the conclusion in the last case.
\end{itemize}\medskip
\ifdefined\hypertarget\hypertarget{sourcepointer-1}{}\fi
\noindent\textit{Source reference.} Original reference decent-spaces-lemma-locally-separated-decent: sources/stacks-algebraic-stacks-morphisms/pinned-original-source.tex, source lines 3965--3965. Navigation pointer only: the counted text refers to the Stacks item decent-spaces-lemma-locally-separated-decent, which belongs to another chapter of the project and is not part of this card. The pointer carries the target's own label as its number because this chapter's staged snapshot does not contain the other chapter's tag table; its content is not reproduced here.\par\medskip
\ifdefined\hypertarget\hypertarget{sourcepointer-2}{}\fi
\noindent\textit{Source reference.} Original reference decent-spaces-section-reasonable-decent: sources/stacks-algebraic-stacks-morphisms/pinned-original-source.tex, source lines 3964--3964. Navigation pointer only: the counted text refers to the Stacks item decent-spaces-section-reasonable-decent, which belongs to another chapter of the project and is not part of this card. The pointer carries the target's own label as its number because this chapter's staged snapshot does not contain the other chapter's tag table; its content is not reproduced here.\par\medskip
\ifdefined\hypertarget\hypertarget{sourcepointer-3}{}\fi
\noindent\textit{Source reference.} Original lemma lemma-diagonal-diagonal: sources/stacks-algebraic-stacks-morphisms/pinned-original-source.tex, source lines 1585--1585. Navigation pointer only: the counted text refers to the Stacks item lemma-diagonal-diagonal, which is an existing labeled result of this same chapter and is not part of this card. The pointer carries the target's own label as its number and its definition line as its span; its content is not reproduced here.\par\medskip
\ifdefined\hypertarget\hypertarget{sourcepointer-4}{}\fi
\noindent\textit{Source reference.} Original reference spaces-more-groupoids-lemma-group-scheme-over-field-separated: sources/stacks-algebraic-stacks-morphisms/pinned-original-source.tex, source lines 3967--3967. Navigation pointer only: the counted text refers to the Stacks item spaces-more-groupoids-lemma-group-scheme-over-field-separated, which belongs to another chapter of the project and is not part of this card. The pointer carries the target's own label as its number because this chapter's staged snapshot does not contain the other chapter's tag table; its content is not reproduced here.\par\medskip
\ifdefined\hypertarget\hypertarget{sourcepointer-5}{}\fi
\noindent\textit{Source reference.} Original reference spaces-more-groupoids-proposition-group-space-scheme-over-field: sources/stacks-algebraic-stacks-morphisms/pinned-original-source.tex, source lines 3970--3970. Navigation pointer only: the counted text refers to the Stacks item spaces-more-groupoids-proposition-group-space-scheme-over-field, which belongs to another chapter of the project and is not part of this card. The pointer carries the target's own label as its number because this chapter's staged snapshot does not contain the other chapter's tag table; its content is not reproduced here.\par\medskip
\ifdefined\hypertarget\hypertarget{sourcepointer-6}{}\fi
\noindent\textit{Source reference.} Original reference spaces-over-fields-lemma-locally-quasi-finite-over-field: sources/stacks-algebraic-stacks-morphisms/pinned-original-source.tex, source lines 3974--3974. Navigation pointer only: the counted text refers to the Stacks item spaces-over-fields-lemma-locally-quasi-finite-over-field, which belongs to another chapter of the project and is not part of this card. The pointer carries the target's own label as its number because this chapter's staged snapshot does not contain the other chapter's tag table; its content is not reproduced here.\par\medskip
\begin{lemma}
\label{lemma-automorphism-group-scheme}
In the situation above $G_x$ is a scheme if one of the following
holds
\begin{enumerate}
\item $\Delta : \mathcal{X} \to \mathcal{X} \times \mathcal{X}$
is quasi-separated
\item $\Delta : \mathcal{X} \to \mathcal{X} \times \mathcal{X}$
is locally separated,
\item $\mathcal{X}$ is quasi-DM,
\item $\mathcal{I}_\mathcal{X} \to \mathcal{X}$
is quasi-separated,
\item $\mathcal{I}_\mathcal{X} \to \mathcal{X}$
is locally separated, or
\item $\mathcal{I}_\mathcal{X} \to \mathcal{X}$
is locally quasi-finite.
\end{enumerate}
\end{lemma}

\begin{proof}
Observe that (1) $\Rightarrow$ (4), (2) $\Rightarrow$ (5), and
(3) $\Rightarrow$ (6) by
Lemma \ref{lemma-diagonal-diagonal}.
In case (4) we see that $G_x$ is a quasi-separated
algebraic space and in case (5) we see that $G_x$
is a locally separated algebraic space.
In both cases $G_x$ is a decent algebraic space
(Decent Spaces, Section \ref{decent-spaces-section-reasonable-decent} and
Lemma \ref{decent-spaces-lemma-locally-separated-decent}).
Then $G_x$ is separated by More on Groupoids in Spaces, Lemma
\ref{spaces-more-groupoids-lemma-group-scheme-over-field-separated}
whereupon we conclude that $G_x$ is a scheme by
More on Groupoids in Spaces, Proposition
\ref{spaces-more-groupoids-proposition-group-space-scheme-over-field}.
In case (6) we see that $G_x \to \Spec(k)$ is locally quasi-finite
and hence $G_x$ is a scheme by
Spaces over Fields, Lemma
\ref{spaces-over-fields-lemma-locally-quasi-finite-over-field}.
\end{proof}

\end{document}
